\chapter{Einleitung}
\label{sec:introduction}
Die digitale Signalverarbeitung (Digital Signal Processing, DSP) ist ein wichtiges Anwendungsfeld in der Informatik. Besonderen Stellenwert hat sie
Sie findet unter anderem Anwendung in der Telekommunikation, Audiotechnik,
Sprachverarbeitung und Medizintechnik \parencite{ifeachor2002digital}. In diesen Anwendungsbereichen sind
Laufzeiteffizienz, deterministische Latenz und numerische Präzision von
besonderer Bedeutung.

DSP-Algorithmen werden traditionell in C oder C++ implementiert, unter
unter anderem wegen des direkten Hardwarezugriffs und der stark optimierten Codeerzeugung \parencite{levine1997dsp101}.

Die manuelle Übertragung mathematischer
Modelle in ausführbaren Code ist jedoch mit erheblichem Aufwand verbunden und erschwert die Portierung
auf unterschiedliche Zielplattformen. In den letzten zwanzig Jahren wurden daher domänenspezifische Sprachen (Domain-Specific Languages, DSLs) für die Audio-Signalverarbeitung entwickelt, die durch Abstraktion und Code-Generierung diese Herausforderungen adressieren sollen. \\parencite[S. 2-6]{orlarey:hal-02159014}

\textcite{orlarey:hal-02159014} beschreibt die domänenspezifische Sprache FAUST (Functional Audio Stream), die am GRAME in Lyon entwickelt wurde, als eine funktionale, blockbasierte Programmiersprache. Sie wurde speziell für Audio-DSP entwickelt. Programme bestehen aus einer algebraischen Kombination von Signalverarbeitungsoperatoren. Die aktuelle FAUST-Dokumentation beschreibt die Codegenerierung für mehrere
Zielsprachen und Zielplattformen, darunter C++, Rust, JavaScript und
WebAssembly \parencite{faust_manual}. Sie dokumentiert zudem
Architekturdateien, mit denen aus einer FAUST-Quellbeschreibung unter anderem
Audio-Plug-ins, Anwendungen für eingebettete Systeme und Webanwendungen
erzeugt werden können \parencite{faust_manual}.

Die praktische Leistungsfähigkeit dieses Ansatzes wird am Beispiel des
FAUST-STK untersucht. FAUST-STK ist eine Sammlung virtueller Musikinstrumente,
deren Modelle überwiegend auf Wellenleitersynthese und teilweise auf modaler
Synthese beruhen \parencite[S.~1]{michon2011fauststk}. Die aus dem Synthesis
ToolKit und SynthBuilder übernommenen Modelle wurden für die FAUST-Semantik
angepasst und hinsichtlich ihrer Effizienz optimiert
\parencite[S.~5--6]{michon2011fauststk}.

Die Quelltextgrößenanalyse zeigt, dass die FAUST-Implementierungen bei den
untersuchten Modellen kompakter sind als die entsprechenden STK-C++-
Implementierungen. Der ausgewiesene Größenvorteil beträgt je nach Modell
zwischen 22,91 Prozent und 80,20 Prozent
\parencite[Tab.~1, S.~6]{michon2011fauststk}. 
Zusätzlich vergleicht die Studie
die CPU-Last von Pure-Data-Plug-ins, die auf STK-C++-Code beziehungsweise
FAUST-generiertem Code basieren \parencite[Tab.~2, S.~6]{michon2011fauststk}.

Für die praktische Eignung ist es erforderlich, dass der durch FAUST generierte Code sowohl funktional korrekt als auch laufzeiteffizient ist.
\textcite{scaringella:hal-02158949} thematisieren die automatische Vektorisierung, während \textcite{letz:hal-02158924} Möglichkeiten der Parallelisierung im FAUST-Umfeld behandeln. 
Die FAUST-Dokumentation verweist zudem auf Optimierungen, die den vollständigen Signalfluss berücksichtigen \parencite{orlarey:hal-02159014}. Diese Eigenschaften lassen jedoch keine allgemeine Aussage darüber zu, ob FAUST-generierter C++-Code einer manuell implementierten C++-Referenz bei konkreten DSP-Bausteinen überlegen, gleichwertig oder unterlegen ist.

Direkte Vergleiche liegen für zeitbereichsbasierte Algorithmen vor. Für den Reverb-Algorithmus \emph{Freeverb} berichten \textcite{orlarey:hal-02159014} Benchmarkergebnisse, deren Ausprägung von Compiler und
Codegenerierungsstrategie abhängt. Die in der vorliegenden Arbeit berücksichtigte Vergleichsliteratur untersucht überwiegend zeitbereichsbasierte Algorithmen wie Reverbs, Delays und physikalische Modelle. Ergänzend wird daher ein Vergleich FFT-bezogener Bausteine durchgeführt.
Für einen fairen Vergleich folgen beide Implementierungen demselben algorithmischen Ablauf und werden über eine gemeinsame Schnittstelle in die Benchmark-Infrastruktur eingebunden. Auf explizit programmierte Single-Instruction-Multiple-Data-Intrinsics (SIMD-Intrinsics) wird verzichtet.
Die automatische Vektorisierung des Compilers bleibt in beiden Varianten aktiv.

\section{Forschungsfrage und Zielsetzung}
\label{sec:forschungsfrage}
Diese Arbeit untersucht das Verhalten von FAUST-generiertem C++-Code im Vergleich zu einer manuell implementierten C++-Referenz für grundlegende DSP-Bausteine. Berücksichtigt werden eine Gain-Stufe, ein Biquad-Filter, Finite-Impulse-Response-Filter (FIR-Filter) verschiedener Längen sowie FFTs unterschiedlicher Größen. Die Anforderungen umfassen reproduzierbare Vergleiche unter unterschiedlichen Datenabhängigkeiten und Skalierungseigenschaften. Neben der Laufzeit werden auch Speicherverbrauch und Binärgröße analysiert. Zusätzlich erfolgt eine Analyse der Codelänge und Implementierungskomplexität.

Die übergeordnete Forschungsfrage lautet:
\begin{quote}
	Wie unterscheidet sich FAUST-generierter C++-Code von einer manuell
	implementierten C++-Referenz bei ausgewählten DSP-Bausteinen und FFTs
	hinsichtlich Laufzeit, Speicherverbrauch, Binärgröße, Codelänge und
	Implementierungskomplexität?
\end{quote}

Zur Beantwortung dieser Forschungsfrage werden folgende Teilfragen untersucht:

\begin{itemize}
	\item Wie groß ist der Performanzunterschied hinsichtlich der CPU-Zeit pro Audio-Block zwischen FAUST-generiertem und manuell implementiertem C++-Code bei Gain-Stufen, Biquad-Filtern, FIR-Filtern und FFTs?
	\item Wie verändert sich der Laufzeitunterschied zwischen beiden
	Implementierungen in Abhängigkeit von Blockgröße, FIR-Filterlänge
	und FFT-Größe?
	\item Wie unterscheiden sich FAUST-generierter und manuell implementierter C++-Code hinsichtlich des Speicherverbrauchs während der Ausführung der untersuchten DSP-Bausteine?
	\item Welche Unterschiede bestehen zwischen FAUST-generiertem und manuell implementiertem C++-Code hinsichtlich der Binärgröße der erzeugten Programme?
	\item Welchen Einfluss haben Compileroptimierungen und automatische Vektorisierung auf beide Implementierungen?
	\item Welche Unterschiede bestehen hinsichtlich Codelänge und Implementierungskomplexität der beiden Ansätze?
\end{itemize}


\section{Aufbau der Arbeit}
Die Bachelorarbeit gliedert sich in acht Kapitel. Sie behandelt die theoretischen Grundlagen, die Beschreibung von FAUST, das Versuchsdesign, die Implementierung der Benchmark-Kerne, die Evaluation und die Einordnung der Ergebnisse.

Nach dieser Einleitung vermittelt Kapitel~\ref{cha:grundlagen} die für die Arbeit relevanten Grundlagen der digitalen Signalverarbeitung. Dazu gehören Audio-DSP, Verstärkungsstufen, rekursive Filter und FIR-Filter, die Fast-Fourier-Transformation sowie Anforderungen an Echtzeitverarbeitung.

Kapitel~\ref{cha:faust} stellt die domänenspezifische Sprache FAUST vor. Es erläutert die funktionale Syntax und Semantik, die zugrunde liegende Compiler-Architektur sowie die für die Arbeit relevanten Codegenerierungs-Backends.


Kapitel~\ref{chap:versuchsdesign} beschreibt das Versuchsdesign und die Benchmark-Methodik. Es definiert die untersuchten DSP-Bausteine, die Messgrößen, die Zielplattformen, die Compilerkonfigurationen und die Maßnahmen zur Sicherung reproduzierbarer Messungen.


Kapitel~\ref{cha:implementierung} dokumentiert die praktische Umsetzung der FAUST- und C++-Varianten. Neben den DSP-Bausteinen werden die gemeinsame Schnittstelle, das Build-System und die Maßnahmen zur Gewährleistung einer fairen Vergleichsbasis beschrieben.


Kapitel~\ref{cha:evaluation} enthält die Evaluation. Die Ergebnisse der Mikrobenchmarks werden hinsichtlich Laufzeit, Speicherverbrauch und Binärgröße ausgewertet. Zusätzlich werden Skalierungseffekte durch unterschiedliche Blockgrößen, Filterlängen und FFT-Größen sowie der Einfluss von Compileroptimierungen analysiert. Im Anschluss wird die Erfahrung mit der Handhabung von FAUST bewertet.


Kapitel~\ref{cha:diskussion} diskutiert die Ergebnisse im Hinblick auf die Forschungsfrage. Dabei werden Limitationen des Versuchsaufbaus, Bedrohungen der Validität und sowie die Codelänge und
Implementierungskomplexität analysiert.

Kapitel~\ref{cha:zusammenfassung} fasst die zentralen Ergebnisse zusammen und gibt einen Ausblick auf mögliche Erweiterungen, etwa die Untersuchung vollständiger DSP-Pipelines oder weiterer Zielplattformen.
